\section{研究变化本身：先识别主导力量}

AI 进展太快时，逐条追新闻会积累事实，却不一定形成方向判断。Hyung Won Chung 选择另一条路线：把 Transformer 的架构史当成一组自然实验，追问哪些设计在旧任务与旧算力下合理，哪些设计随着规模、数据和交互方式变化而失去优势。整堂课的目标不是证明 decoder-only 永远正确，而是训练一种可迁移的研究方法。

\subsection{为什么历史分析不是怀旧}

标题页把 ``history'' 与 ``future'' 放在同一句话里。读下面两页时应先抓住方法论：研究对象不是某个模型名称，而是变化轨迹。只要能识别支配轨迹的方向性变量，就能把杂乱更新压缩成少数可检验假设。后文每一次架构比较都要回到这个问题：它是在记录一个静态赢家，还是在解释为什么相对优势会随 regime 移动？

\lecturefigure{hyung-slide-001.jpg}{讲座主题：从 Transformer 历史塑造 AI 的未来}{Hyung Won Chung official deck, p.~1}

\lecturefigure{hyung-slide-002.jpg}{不要只追最新结果，要研究变化本身}{Hyung Won Chung official deck, p.~2}

\teachervoice{00:02:05--00:02:58，Hyung Won 强调：当变化速度超过人的追踪能力时，应回看旧系统，画出“我们怎样走到现在”的路径。课堂重点不是背 2024 年的榜单，而是从路径中抽象方向。}

\begin{warningbox}{独立剪辑的证据边界}
本视频是官方播放列表中的独立 36:30 剪辑，与 Lecture 27 合并录像的 Hyung Won 后半段内容相同，但不包含后续联合 Q\&A。本讲义只保留独立视频中实际出现的 prepared talk，不把 Q\&A 中关于 MLE、RLHF 或能源极限的讨论倒灌进来。
\end{warningbox}

\subsection{Identify、understand、roll forward}

第三页把预测拆成三步：先识别 dominant driving force，再理解它怎样改变方法之间的相对优劣，最后谨慎地把趋势向前外推。主导力量不是唯一力量，而是在当前精度要求下解释大部分方向性的变量；如果 regime 改变，就必须重新建模。读图前还应区分“描述过去”与“预测未来”：前者可由历史数据约束，后者必须同时写出失效条件和替代解释。

\lecturefigure{hyung-slide-003.jpg}{研究变化的三步法：识别、理解并外推主导力量}{Hyung Won Chung official deck, p.~3}

\begin{importantbox}{Dominant driving force 的工作定义}
主导驱动力是在给定时间尺度、预算与误差容忍度下，足以解释主要方向变化的因素。它允许研究者暂时忽略较小作用，以换取可分析性；但忽略项必须公开，预测失败时也应首先检查被省略的反馈与约束。
\end{importantbox}

\teachervoice{00:03:03--00:04:07，讲者把研究流程概括为 identify、understand、roll forward，并指出预测命中率不必很高才有价值。重要的是让假设可检查，而不是假装精确知道未来。}

\subsection{落笔实验：什么时候可以忽略空气阻力}

一支笔从静止释放时，现实中同时存在重力、空气阻力、旋转与材料形变；课堂仍只建模重力，因为在短距离、普通精度下它足够主导。若把向下方向记为正，可写为
\[
y(t)=y_0+v_0t+\frac{1}{2}gt^2.
\]
其中 $y(t)$ 是时刻 $t$ 的位置，$y_0$ 是初始位置，$v_0$ 是初速度，$g$ 是重力加速度。公式成立的条件不是“阻力不存在”，而是“阻力对当前问题的误差贡献可接受”。

\lecturefigure{hyung-slide-004.jpg}{落笔实验：用重力说明如何忽略次要力量}{Hyung Won Chung official deck, p.~4}

\lecturefigure{hyung-slide-006.jpg}{主导力量越多、相互作用越复杂，轨迹越难预测}{Hyung Won Chung official deck, p.~6}

\begin{knowledgebox}{读图：横轴不是观测字段数量}
横轴表示\textbf{相互竞争的主导力量数量}。一个系统可以有很多特征，却仍由少数机制决定方向；反过来，变量很少但反馈强耦合，也可能难以外推。图支持的是建模顺序：先找方向性最强的因素，再决定何时加入相互作用。
\end{knowledgebox}

\teachervoice{00:04:40--00:05:57，讲者明确提到空气阻力和复杂气动相互作用，却选择忽略它们。课堂提示不是“现实只有一个变量”，而是“先定义精度需求，再决定哪些变量可以近似掉”。}

\subsection{本章小结}

本章建立了整堂课的证据纪律：历史不是模型名单，而是自然实验；主导力量不是唯一因果，而是可声明误差边界的近似；外推不是保证，而是生成下一组实验。接下来把这套方法用于 compute、inductive bias 与 Transformer 架构。

